% !TEX root = IE3_expA_report.tex
\documentclass[lualatex,a4paper,10pt]{jlreq}
% Shared LuaLaTeX preamble.
\usepackage{luatexja-fontspec}
\setmainfont{TeX Gyre Termes}
\setsansfont{TeX Gyre Heros}
\setmonofont{Latin Modern Mono}
\setmainjfont[BoldFont=HaranoAjiGothic-Medium]{HaranoAjiMincho-Regular}
\setsansjfont[BoldFont=HaranoAjiGothic-Bold]{HaranoAjiGothic-Medium}
\usepackage[top=22mm,bottom=22mm,left=23mm,right=23mm]{geometry}
\usepackage{amsmath,amssymb,booktabs,array,longtable,tabularx}
\usepackage{graphicx,xcolor,tikz,circuitikz}
\usetikzlibrary{arrows.meta,positioning,calc}
\usepackage{enumitem,fancyhdr,listings,needspace}
\usepackage[unicode,colorlinks=true,linkcolor=labblue,urlcolor=labteal]{hyperref}
\usepackage{bookmark}
\definecolor{labblue}{RGB}{30,70,112}
\definecolor{labteal}{RGB}{0,100,105}
\definecolor{labgray}{gray}{0.95}
\ModifyHeading{section}{font={\large\sffamily\bfseries\color{labblue}}}
\ModifyHeading{subsection}{font={\normalsize\sffamily\bfseries\color{labteal}}}
\setlength{\parindent}{1\zw}
\setlength{\parskip}{3pt}
\setlength{\emergencystretch}{2em}
\renewcommand{\arraystretch}{1.35}
\setlist{itemsep=3pt,topsep=4pt,leftmargin=2\zw}
\newcolumntype{Y}{>{\raggedright\arraybackslash}X}
\newcolumntype{C}[1]{>{\centering\arraybackslash}p{#1}}
\pagestyle{fancy}
\fancyhf{}
\fancyhead[L]{\small\color{labblue} IE3 後期・電子工学実験}
\fancyhead[R]{\small テーマ\LabCode}
\fancyfoot[C]{\thepage}
\setlength{\headheight}{16pt}
\DeclareOldFontCommand{\rm}{\normalfont\rmfamily}{\mathrm}
\lstset{basicstyle=\small\ttfamily,columns=fullflexible,keepspaces=true,
 breaklines=true,frame=single,backgroundcolor=\color{labgray},
 numbers=left,numberstyle=\tiny,showstringspaces=false,aboveskip=8pt,belowskip=8pt}
\newcommand{\blank}{\rule{22mm}{0.3pt}}
\newcommand{\answerline}{\par\noindent\rule{\linewidth}{0.25pt}\par}
\newcommand{\writing}[1]{\par\Needspace{\dimexpr#1+5mm\relax}\noindent%
 \fcolorbox{labblue!35}{white}{\parbox[c][#1][t]{\dimexpr\linewidth-2\fboxsep-2\fboxrule\relax}{\vspace{2mm}\noindent\strut\hfill\par}}\par}
\newcommand{\hint}[1]{\par\smallskip\noindent{\sffamily\bfseries\color{labteal}記述の視点：}#1\par}
\newcommand{\note}[1]{\par\smallskip\noindent\fcolorbox{labblue!45}{labblue!5}{%
 \parbox{\dimexpr\linewidth-2\fboxsep-2\fboxrule\relax}{#1}}\par\smallskip}
\newcommand{\eqerror}{\varepsilon=\frac{x_{\mathrm{meas}}-x_{\mathrm{th}}}{x_{\mathrm{th}}}\times100\ \%}
\newcommand{\LabSource}{徳山工業高等専門学校 情報電子工学科，
 『2026年度 電子工学実験（3年後期）テキスト』，テーマ\LabCode，
 \expandafter\path\expandafter{\SourcePdf}。}
\newcommand{\reporttitle}{
 \noindent{\small\sffamily\color{labteal}実験報告書・記入ガイド}\par
 \noindent{\LARGE\sffamily\bfseries \LabCode\quad\LabTitle}\par\smallskip
 \noindent 氏名：\blank\quad 班：\blank\quad 実験日：\blank\par
 \note{目的・原理・手順を確認し、実測値と自分の考察を記入する。
 考察のガイドは、検討する事象・使う測定結果・記述の方針を示す。}
}
\newcommand{\prelabtitle}{
 \noindent{\small\sffamily\color{labteal}事前学習・前提知識プリント}\par
 \noindent{\LARGE\sffamily\bfseries \LabCode\quad\LabTitle}\par\smallskip
 \noindent 氏名：\blank\quad 班：\blank\quad 予習日：\blank\par
 \note{用語、回路の動き、計算、測定表への記入の順に読む。
 例題の値は説明用であり、実験結果ではない。}
}
\newcommand{\tocrow}[3]{\hyperref[#1]{#2}& #3 &\pageref*{#1}\\[2mm]}
\newcommand{\documentmap}[1]{
 \section*{目次と概要}
 \noindent 青色の項目をクリックすると該当箇所へ移動する。\par
 {\small\begin{longtable}{@{}p{0.29\linewidth}p{0.61\linewidth}r@{}}
 \toprule {\color{labblue}\bfseries 項目} & {\color{labblue}\bfseries 読むと分かること} & 頁\\\midrule
 \endhead #1\bottomrule
 \end{longtable}}
 \clearpage
}
\newenvironment{terms}{\par\smallskip\begin{longtable}{@{}p{0.23\linewidth}p{0.73\linewidth}@{}}
 \toprule {\color{labteal}\bfseries 用語}&{\color{labteal}\bfseries 意味・回路での役割}\\\midrule\endhead}
 {\bottomrule\end{longtable}}
\newcommand{\term}[2]{\textbf{#1}&#2\\[2mm]}
\newcommand{\guidepoint}[4]{
 \Needspace{32mm}\subsection{#1}
 \noindent{\bfseries\color{labteal}考える事象：}#2\par
 \noindent{\bfseries\color{labteal}参照する結果：}#3\par
 \noindent{\bfseries\color{labteal}記述の方針：}#4\par
}
\newcommand{\reportclosing}[1]{
 \clearpage\section{結論のまとめ方}\label{sec:closing}
 \hint{#1}
 \ConclusionGuide
 \begin{enumerate}
 \item 目的に対応する最も重要な実測結果を、測定条件と数値を添えて述べる。
 \item 原理と一致した点・一致しなかった点を示す。原因は確認済みの事実と推測を分ける。
 \item 実施範囲で言えることと、未確認の条件を一文ずつまとめる。
 \end{enumerate}
 \note{書き出しの型：「○○の条件で△△を測定したところ、□□となった。
 これは◇◇と整合する。ただし××は確認しておらず、一般化には追加測定が必要である。」
 空欄は自分の実測値と根拠で埋める。}
 \noindent 結論（実験後に記入）：\writing{30mm}
 \noindent 改善案・次に確認したい条件：\writing{16mm}
 \section*{参考文献}\label{sec:references}
 \begin{enumerate}
 \item \LabSource
 \item 使用したデータシート・書籍・ウェブ資料（著者、題名、該当箇所、URL、参照日）を追加する。
 \end{enumerate}
}
\newcommand{\prelabclosing}{
 \section*{準備・出典}\label{sec:prelabsource}
 \begin{itemize}[nosep]
 \item 資料、筆記具、記録用紙、電卓と、実験条件・機器設定の記録欄。
 \item 確認問題の解答と、分からない用語・回路点への具体的な質問。
 \end{itemize}
 {\small\noindent\LabSource\par}
}
\newcommand{\simpleflow}[1]{
 \begin{center}
 \begin{tikzpicture}[node distance=5mm and 5mm,
 every node/.style={font=\small},box/.style={draw=labblue,fill=labblue!4,rounded corners=1mm,
 align=center,minimum height=10mm,text width=30mm}]
 #1
 \end{tikzpicture}
 \end{center}
}

\newcommand{\LabCode}{A}
\newcommand{\LabTitle}{回路網}
\newcommand{\SourcePdf}{IE3_expA_A4.pdf}
\newcommand{\ConclusionGuide}{三つの定理を一文ずつ対応させる。テブナンは「どの抵抗を含めた予測か」、重ねの理は「符号付き和と残差」、相反定理は「交換前後の差」を根拠として使う。
「成立した」という表現を使うときは、測定した電源条件と読取り精度の範囲を添える。\par}
\begin{document}
\reporttitle
\documentmap{
\tocrow{sec:auto1}{1 目的}{端子から見た等価回路と回路網定理を実験で確かめる。}
\tocrow{sec:auto2}{2 原理}{端子から見た等価回路と回路網定理の仕組みと成立条件を整理する。}
\tocrow{sec:auto3}{3 使用器具}{機器・定数・測定点を確認し、資料の順序で操作する。}
\tocrow{sec:auto4}{4 実験方法}{機器・定数・測定点を確認し、資料の順序で操作する。}
\tocrow{sec:auto5}{5 結果}{電流の予測、符号付き和、交換前後の差を表や図に記録する。}
\tocrow{sec:auto6}{6 考察：結果を根拠に説明する}{電流の予測、符号付き和、交換前後の差を根拠に、比較と原因の検討を進める。}
\tocrow{sec:closing}{7 結論のまとめ方}{主要な結果、成立範囲、未確認の条件を短くまとめる。}
\tocrow{sec:references}{参考文献}{使用した資料と外部の根拠を示す。}
}


\section{目的}\label{sec:auto1}
鳳・テブナンの定理、重ねの理、相反定理を実測により確かめ、各定理の適用条件と測定器の影響を理解する。
\section{原理}\label{sec:auto2}
\subsection{鳳・テブナンの定理}
線形回路網を端子から見ると、開放端子電圧 $V_{\rm th}$ の電圧源と等価抵抗 $R_{\rm th}$ の直列回路に置き換えられる。
回路網Bの抵抗を $R_B$、電流計の内部抵抗を $r_A$ とすると、
\begin{equation} I_{\rm th}=\frac{V_{\rm th}}{R_{\rm th}+R_B+r_A}.\label{eq:th}\end{equation}
独立電圧源の作用を除くときは、電源を外してその接続点を短絡する。通電中の電源出力を直接短絡する操作とは区別する。
\subsection{重ねの理と相反定理}
複数電源がある線形回路では、同じ基準方向の電流・電圧について
\begin{equation}I=I_1+I_2,\qquad V=V_1+V_2\end{equation}
が成り立つ。理想電圧源をゼロにすると短絡、理想電流源をゼロにすると開放になる。
線形かつ双方向性をもつ抵抗回路では、電源位置と電流観測位置を交換したときの伝達関係が等しい。
資料A-2.4の補償定理は抵抗変化による電流変化を等価起電力で扱う考え方であり、今回の実測項目はA-3.1〜A-3.3である。
\section{使用器具}\label{sec:auto3}
直流電源2台、直流電圧計、直流電流計、可変抵抗器、検流計、ホイートストンブリッジ用回路、回路網A・Bを使用する。
\par\noindent\begin{tabularx}{\linewidth}{YYY}\toprule 器具 & 型番・個体番号 & レンジ・備考\\\midrule
直流電源1／2&&\\電圧計／電流計&&\\抵抗測定回路・検流計&&\\回路網A／B&&\\\bottomrule\end{tabularx}
\clearpage
\section{実験方法}\label{sec:auto4}
\subsection{鳳・テブナンの定理（資料A-3.1）}
\begin{enumerate}
\item 資料図A-6に従い、回路網Aの電源端子に $E=2,4,6,8,10$ Vを順に加え、負荷未接続時の電圧を測定する。
\item 電源を外して電源接続点を短絡し、ブリッジで $R_{\rm th}$ を測定する。回路網Bの $R_B$ も測定する。ブリッジの電源は2〜3 Vとする。
\item 図A-7のようにA・Bと電流計を接続し、同じ5条件で電流を測定する。電流計のレンジと内部抵抗は実機表示から記録する。
\end{enumerate}
\subsection{重ねの理（資料A-3.2）}
$(E_1,E_2)=(2,5),(4,4),(6,3),(8,2),(10,1)$ Vを設定し、両電源接続時の $V,I$、
$E_1$のみの $V_1,I_1$、$E_2$のみの $V_2,I_2$ を測定する。
取り外した電源の接続点は短絡し、電圧極性・電流方向を統一する。
\subsection{相反定理（資料A-3.3）}
図A-9の2接続について、$E=2,4,6,8,10$ Vで電流 $I_a,I_b$ を測定する。
電源と電流計を入れ替え、計器レンジを維持する。端子の緩みを確認する。
\section{結果}\label{sec:auto5}
\noindent $R_{\rm th}=$\blank $\Omega$，$R_B=$\blank $\Omega$，
$r_A=$\blank $\Omega$，電流計レンジ：\blank
\begin{table}[h]\centering\caption{テブナンの定理の測定と予測}
\begin{tabular}{ccccc}\toprule $E$ [V]&$V_{\rm th}$ [V]&$I$ [mA]&$I_{\rm th}$ [mA]&誤差 [\%]\\\midrule
2&&&&\\4&&&&\\6&&&&\\8&&&&\\10&&&&\\\bottomrule\end{tabular}\end{table}
式\eqref{eq:th}では電流をAで計算した後にmAへ換算する。
\clearpage
\begin{table}[h]\centering\caption{重ねの理：符号を含む測定値}
\begin{tabular}{cccccccc}\toprule $E_1$&$E_2$&$V$&$I$&$V_1$&$I_1$&$V_2$&$I_2$\\
\relax [V]&[V]&[V]&[mA]&[V]&[mA]&[V]&[mA]\\\midrule
2&5&&&&&&\\4&4&&&&&&\\6&3&&&&&&\\8&2&&&&&&\\10&1&&&&&&\\\bottomrule\end{tabular}\end{table}
\noindent 残差 $\Delta V=V-(V_1+V_2)$，$\Delta I=I-(I_1+I_2)$：
\writing{15mm}
\begin{table}[h]\centering\caption{相反定理の確認}
\begin{tabular}{cccc}\toprule $E$ [V]&$I_a$ [mA]&$I_b$ [mA]&$I_a-I_b$ [mA]\\\midrule
2&&&\\4&&&\\6&&&\\8&&&\\10&&&\\\bottomrule\end{tabular}\end{table}
\noindent 実際の接続図（資料図A-7、A-8、A-9を参考に測定方向を明示）：
\writing{25mm}
\clearpage\section{考察：結果を根拠に説明する}\label{sec:auto6}
\guidepoint{テブナン予測と計器の影響}
{予測電流と実測電流の差が、電流計の内部抵抗を含めると変わるか。}
{表1の$V_{\rm th},I$、測定した$R_{\rm th},R_B,r_A$、使用レンジ。}
{各条件で内部抵抗あり・なしの二予測を計算する。どちらが実測に近いかを数値で示す。差が残る場合は抵抗公差や電圧計負荷を候補とし、未測定の原因は推測として述べる。}
\guidepoint{重ね合わせの一致をどう判断するか}
{二電源同時の値が、一電源ずつの符号付き和とどの程度一致するか。}
{表2の$V,I,V_1,I_1,V_2,I_2$、残差、正方向を描いた接続図。}
{電圧と電流を別々に比較する。残差を単位付きで示し、零付近では相対誤差より絶対差と計器の読取り幅を使う。極性・短絡操作の取り違えも検討する。}
\guidepoint{相反定理の交換条件}
{励振位置と測定位置の交換で電流が変わるか。}
{表3の同じEにおける$I_a,I_b$、二つの接続図、計器レンジ。}
{$I_a-I_b$を各行で求める。電源・レンジが同じであることを確認し、一致した範囲を述べる。双方向性を持たない素子への一般化は行わない。}
\guidepoint{ばらつきと再現性}
{理論との差が測定の散らばりより大きいか。}
{再測定値、端子の状態、電源設定、最小目盛の記録。}
{再測定があれば範囲を示す。なければ再現性は未確認とし、同じ条件で繰り返す計画を改善案として挙げる。}
\subsection{資料の課題・追加検討}
\begin{enumerate}
\item 内部抵抗を含めた予測と含めない予測を比較し、計器の負荷効果を定量的に述べる。
\item 重ねの理の残差を読取り誤差や抵抗・電源の変動と比較する。ゼロ付近では相対誤差だけで判断しない。
\item 相反定理の成立を計器条件と測定の不確かさを含めて判断する。
\end{enumerate}
\writing{12mm}
\reportclosing{3つの定理について、どの程度実測と一致したかをまとめる。}

\end{document}
